\documentclass[border=10pt]{standalone}
\usepackage{ctex}
\usepackage[european]{circuitikz}
\begin{document}
\begin{circuitikz}
% 测试 dipchip 的自定义引脚标注
\ctikzset{multipoles/dipchip/width=2.0}
\draw (0,0) node[dipchip, num pins=12, hide numbers, no marks] (c) {};
\node[right, font=\small] at (c.bpin 1) {CP$_0$};
\node[right, font=\small] at (c.bpin 2) {CP$_1$};
\node[right, font=\small] at (c.bpin 3) {R$_0$(1)};
\node[right, font=\small] at (c.bpin 4) {R$_0$(2)};
\node[right, font=\small] at (c.bpin 5) {S$_9$(1)};
\node[right, font=\small] at (c.bpin 6) {S$_9$(2)};
\node[left, font=\small] at (c.bpin 12) {Q$_3$};
\node[left, font=\small] at (c.bpin 11) {Q$_2$};
\node[left, font=\small] at (c.bpin 10) {Q$_1$};
\node[left, font=\small] at (c.bpin 9) {Q$_0$};

% 测试 dipchip 顶部标签
\node[above] at (c.north) {74LS290};

% 测试 74LS161
\draw (5,0) node[dipchip, num pins=12, hide numbers, no marks] (c2) {};
\node[above] at (c2.north) {74LS161};
\node[right, font=\small] at (c2.bpin 1) {$\overline{\text{R}_\text{D}}$};
\node[right, font=\small] at (c2.bpin 2) {$\overline{\text{L}_\text{D}}$};
\node[right, font=\small] at (c2.bpin 3) {D$_3$D$_2$D$_1$D$_0$};
\node[right, font=\small] at (c2.bpin 4) {CP};
\node[right, font=\small] at (c2.bpin 5) {EP ET};
\node[left, font=\small] at (c2.bpin 9) {CO};
\node[left, font=\small] at (c2.bpin 10) {Q$_3$Q$_2$Q$_1$Q$_0$};

% 测试 european nand port 与 dipchip 的组合
\draw (5,3.5) node[european nand port] (nand) {};
\end{circuitikz}
\end{document}
